\documentclass[12pt,a4paper]{article}

\usepackage[utf8]{inputenc}
\usepackage[T1]{fontenc}
\usepackage[french]{babel}
\usepackage{geometry}
\usepackage{graphicx}
\usepackage{xcolor}
\usepackage{hyperref}
\usepackage{listings}
\usepackage{array}
\usepackage{longtable}

\geometry{margin=2.5cm}

\hypersetup{
    colorlinks=true,
    linkcolor=blue,
    urlcolor=blue
}

\lstset{
    basicstyle=\ttfamily\small,
    breaklines=true,
    frame=single,
    backgroundcolor=\color{gray!8},
    columns=fullflexible
}

\title{Compte rendu -- Catalog Service API}
\author{Yannis GONOS}
\date{\today}

\begin{document}

\maketitle

\tableofcontents
\newpage

\section{Présentation du projet}

Le projet \textbf{Catalog Service} est un micro-service développé avec \textbf{Django} et \textbf{Django REST Framework}. Il permet de gérer un catalogue de produits dans le cadre du projet FlexShop.

L'API expose des endpoints REST permettant de créer, consulter, modifier, supprimer, filtrer, rechercher et trier des produits. Une documentation OpenAPI est également fournie via Swagger UI.

\section{Lien vers le dépôt Git}

Le dépôt Git du projet est disponible à l'adresse suivante :

\begin{center}
\url{A_COMPLETER_AVEC_LE_LIEN_DU_DEPOT_GIT}
\end{center}

\textbf{À compléter avant rendu :} remplacer \texttt{A\_COMPLETER\_AVEC\_LE\_LIEN\_DU\_DEPOT\_GIT} par le lien réel du dépôt Git.

\section{Technologies utilisées}

\begin{itemize}
    \item Python 3.12+
    \item Django 6.0.5
    \item Django REST Framework 3.17.1
    \item django-filter
    \item drf-spectacular pour la documentation OpenAPI
    \item SQLite pour la base de données locale de développement
\end{itemize}

\section{Fonctionnalités réalisées}

\subsection{CRUD complet}

L'API permet d'effectuer toutes les opérations CRUD sur les produits :

\begin{longtable}{|p{4cm}|p{8cm}|}
\hline
\textbf{Méthode et endpoint} & \textbf{Description} \\
\hline
\texttt{GET /api/v1/products/} & Liste tous les produits avec pagination. \\
\hline
\texttt{POST /api/v1/products/} & Crée un nouveau produit. \\
\hline
\texttt{GET /api/v1/products/\{id\}/} & Récupère le détail d'un produit. \\
\hline
\texttt{PATCH /api/v1/products/\{id\}/} & Modifie partiellement un produit. \\
\hline
\texttt{PUT /api/v1/products/\{id\}/} & Remplace complètement un produit. \\
\hline
\texttt{DELETE /api/v1/products/\{id\}/} & Supprime un produit. \\
\hline
\end{longtable}

\subsection{Pagination, filtrage, recherche et tri}

La pagination est configurée globalement dans Django REST Framework avec une taille de page de 10 éléments.

L'API permet également :

\begin{itemize}
    \item le filtrage par catégorie avec \texttt{?category=electronics} ;
    \item la recherche par nom ou description avec \texttt{?search=souris} ;
    \item le tri par prix, date de création ou stock avec \texttt{?ordering=price\_cents}.
\end{itemize}

\subsection{Endpoint custom avec \texttt{@action}}

Un endpoint personnalisé a été ajouté avec le décorateur \texttt{@action} de Django REST Framework :

\begin{lstlisting}
GET /api/v1/products/low-stock/?threshold=5
\end{lstlisting}

Cet endpoint retourne les produits dont le stock est inférieur ou égal au seuil indiqué. Il permet par exemple d'identifier rapidement les produits à réapprovisionner.

\subsection{Validation métier customisée}

Le serializer applique plusieurs règles métier :

\begin{itemize}
    \item le prix doit être strictement positif ;
    \item le prix ne peut pas dépasser \texttt{99 999 999,99} euros ;
    \item le stock ne peut pas dépasser \texttt{100 000} unités ;
    \item le paramètre \texttt{threshold} de l'endpoint \texttt{low-stock} doit être un entier.
\end{itemize}

\subsection{Prix stocké en centimes}

Le prix est stocké en base de données sous forme d'entier avec le champ \texttt{price\_cents}. Cette approche évite les erreurs d'arrondi liées aux nombres flottants et correspond à une pratique courante dans les applications financières.

L'API reste simple à utiliser : l'utilisateur envoie un prix au format décimal, par exemple \texttt{29.99}, et le serializer le convertit automatiquement en \texttt{2999} centimes.

\section{Conformité REST}

Les URI respectent les bonnes pratiques REST :

\begin{itemize}
    \item les ressources sont nommées au pluriel : \texttt{/products/} ;
    \item les URI ne contiennent pas de verbe ;
    \item les méthodes HTTP indiquent l'action à effectuer ;
    \item les codes de statut HTTP sont cohérents : \texttt{200}, \texttt{201}, \texttt{204}, \texttt{400}, \texttt{404}.
\end{itemize}

\section{Documentation OpenAPI et Swagger UI}

La documentation OpenAPI est générée avec \texttt{drf-spectacular}. Elle est disponible via les endpoints suivants :

\begin{itemize}
    \item Swagger UI : \url{http://localhost:8000/api/docs/}
    \item Schéma OpenAPI : \url{http://localhost:8000/api/schema/}
\end{itemize}

\subsection{Capture d'écran de Swagger UI fonctionnel}

La capture d'écran de Swagger UI doit être placée dans le dossier \texttt{rendu/images/} avec le nom :

\begin{lstlisting}
swagger-ui.png
\end{lstlisting}

Une fois l'image ajoutée, elle apparaîtra ci-dessous dans le PDF.

\begin{figure}[h!]
    \centering
    \IfFileExists{images/swagger-ui.png}{
        \includegraphics[width=0.95\textwidth]{images/swagger-ui.png}
    }{
        \fbox{\parbox{0.9\textwidth}{\centering Capture Swagger UI à ajouter : \texttt{rendu/images/swagger-ui.png}}}
    }
    \caption{Swagger UI fonctionnel pour l'API Catalog Service.}
\end{figure}

\section{Collection Postman}

Une collection Postman exportée au format JSON est fournie avec le rendu :

\begin{lstlisting}
rendu/catalog-service-postman-collection.json
\end{lstlisting}

Cette collection contient les requêtes principales suivantes :

\begin{itemize}
    \item liste des produits ;
    \item création d'un produit ;
    \item récupération du détail d'un produit ;
    \item modification partielle ;
    \item suppression ;
    \item filtrage par catégorie ;
    \item recherche ;
    \item tri par prix en centimes ;
    \item endpoint custom \texttt{low-stock} ;
    \item récupération du schéma OpenAPI.
\end{itemize}

\section{Tests unitaires}

Le projet contient 7 tests unitaires. Ils vérifient notamment :

\begin{itemize}
    \item l'accès à la liste des produits ;
    \item la création d'un produit ;
    \item la conversion du prix en centimes ;
    \item le retour \texttt{404} pour un produit inexistant ;
    \item le filtrage par catégorie ;
    \item le rejet d'un prix négatif ;
    \item le fonctionnement de l'endpoint custom \texttt{low-stock} ;
    \item la mise à jour d'un prix avec recalcul de \texttt{price\_cents}.
\end{itemize}

Les tests ont été exécutés avec succès :

\begin{lstlisting}
Found 7 test(s).
Creating test database for alias 'default'...
System check identified no issues (0 silenced).
.......
----------------------------------------------------------------------
Ran 7 tests in 0.028s

OK
Destroying test database for alias 'default'...
\end{lstlisting}

\section{Installation et exécution}

\subsection{Installation des dépendances}

\begin{lstlisting}
python3 -m venv venv
source venv/bin/activate
pip install -r requirements.txt
\end{lstlisting}

\subsection{Application des migrations}

\begin{lstlisting}
python3 manage.py migrate
\end{lstlisting}

\subsection{Lancement du serveur}

\begin{lstlisting}
python3 manage.py runserver
\end{lstlisting}

\subsection{Lancement des tests}

\begin{lstlisting}
python3 manage.py test products
\end{lstlisting}

\section{Qualité du projet}

Le projet contient les éléments attendus pour un rendu propre :

\begin{itemize}
    \item un fichier \texttt{README.md} clair ;
    \item un fichier \texttt{requirements.txt} listant les dépendances ;
    \item un fichier \texttt{.gitignore} excluant notamment \texttt{venv/}, \texttt{db.sqlite3}, \texttt{\_\_pycache\_\_/} et les fichiers \texttt{*.pyc} ;
    \item des tests unitaires fonctionnels ;
    \item une documentation Swagger accessible.
\end{itemize}

\section{Conclusion}

Le micro-service Catalog Service répond aux critères demandés : endpoints CRUD fonctionnels, respect des conventions REST, documentation OpenAPI, tests unitaires, qualité du projet et fonctionnalités bonus. Les bonus intégrés sont l'endpoint custom avec \texttt{@action}, la validation métier customisée et le stockage du prix en centimes.

\end{document}
